\section{Resultados e discussão}
\label{resultados}

Foram medidas \nCombBorda{} combinações na borda e \nCombNuvem{} na nuvem, com \nQuadrosValidos{} quadros válidos cada, das quais se derivaram \nCenarios{} cenários de taxa da câmera. Os tempos de borda vêm de ARM64 nativo (Neoverse-N1, 2 núcleos) e são usados em comparações relativas.

\subsection{Qualidade da detecção}

A Tabela~\ref{tab:qualidade} mostra o mAP@50 da classe pessoa. O YOLOv8n teve a maior qualidade e foi o único modelo cujo mAP@50 aumentou com a resolução (\mapYoloFpTres{}, \mapYoloFpQuatro{} e \mapYoloFpSeis{} em FP32 a 320, 480 e 640 pixels), porque é o único cuja entrada acompanha a resolução da câmera. No SSD-MobileNet~V1 (\mapSsdFpTres{} a 320) e no EfficientDet-Lite0 (\mapEffdetFpTres{} a 320) a resolução da câmera quase não alterou o resultado, já que o quadro é reduzido para a entrada fixa. A quantização INT8 custou pouca qualidade: no YOLOv8n a 640 o mAP@50 passou de \mapYoloFpSeis{} para \mapYoloIntSeis{}, no EfficientDet-Lite0 a 320 de \mapEffdetFpTres{} para \mapEffdetIntTres{}, e no SSD-MobileNet~V1 a variação ficou em milésimos (\mapSsdFpTres{} e \mapSsdIntTres{} a 320).

% Gerado por src/analysis/artigo_tex.py. Não edite à mão.
\begin{table}[htbp]
\centering
\scriptsize
\setlength{\tabcolsep}{3pt}
\caption{Qualidade da detecção de pessoas (mAP@50 e F1 com limiar 0,5) nas 350 imagens do COCO val2017.}
\label{tab:qualidade}
\begin{tabular}{llllrr}
\hline
Parte & Modelo & Precisão & Resolução & mAP@50 & F1 \\
\hline
borda & YOLOv8n & FP32 & 320 & 0,679 & 0,602 \\
borda & YOLOv8n & FP32 & 480 & 0,750 & 0,668 \\
borda & YOLOv8n & FP32 & 640 & 0,770 & 0,690 \\
borda & YOLOv8n & INT8 & 320 & 0,676 & 0,599 \\
borda & YOLOv8n & INT8 & 480 & 0,744 & 0,645 \\
borda & YOLOv8n & INT8 & 640 & 0,760 & 0,671 \\
borda & SSD-MobileNet V1 & FP32 & 320 & 0,603 & 0,595 \\
borda & SSD-MobileNet V1 & FP32 & 480 & 0,603 & 0,601 \\
borda & SSD-MobileNet V1 & FP32 & 640 & 0,602 & 0,603 \\
borda & SSD-MobileNet V1 & INT8 & 320 & 0,605 & 0,603 \\
borda & SSD-MobileNet V1 & INT8 & 480 & 0,605 & 0,611 \\
borda & SSD-MobileNet V1 & INT8 & 640 & 0,610 & 0,611 \\
borda & EfficientDet-Lite0 & FP32 & 320 & 0,514 & 0,519 \\
borda & EfficientDet-Lite0 & FP32 & 480 & 0,512 & 0,517 \\
borda & EfficientDet-Lite0 & FP32 & 640 & 0,507 & 0,516 \\
borda & EfficientDet-Lite0 & INT8 & 320 & 0,472 & 0,502 \\
borda & EfficientDet-Lite0 & INT8 & 480 & 0,468 & 0,495 \\
borda & EfficientDet-Lite0 & INT8 & 640 & 0,469 & 0,492 \\
nuvem & YOLOv8n & FP32 & 320 & 0,668 & 0,591 \\
nuvem & YOLOv8n & FP32 & 480 & 0,745 & 0,653 \\
nuvem & YOLOv8n & FP32 & 640 & 0,761 & 0,680 \\
nuvem & SSD-MobileNet V1 & FP32 & 320 & 0,592 & 0,555 \\
nuvem & SSD-MobileNet V1 & FP32 & 480 & 0,595 & 0,583 \\
nuvem & SSD-MobileNet V1 & FP32 & 640 & 0,600 & 0,583 \\
nuvem & EfficientDet-Lite0 & FP32 & 320 & 0,506 & 0,511 \\
nuvem & EfficientDet-Lite0 & FP32 & 480 & 0,512 & 0,521 \\
nuvem & EfficientDet-Lite0 & FP32 & 640 & 0,498 & 0,507 \\
\hline
\end{tabular}
\end{table}


\subsection{Desempenho na borda (QP1)}

A Tabela~\ref{tab:latencia_borda} resume a latência fim a fim por quadro. A quantização INT8 reduziu a latência em todas as combinações, com diferença significativa em todos os pares (última coluna da Tabela~\ref{tab:latencia_borda}): o ganho foi de cerca de \razaoYoloTres{} vezes no YOLOv8n, \razaoSsdTres{} no SSD-MobileNet~V1 e \razaoEffdetTres{} no EfficientDet-Lite0 (quadro de 320 pixels). No YOLOv8n a latência cresce com a resolução (\etoeYoloIntTres{}, \etoeYoloIntQuatro{} e \etoeYoloIntSeis{}\,ms em INT8), enquanto nos modelos de entrada fixa só o custo de redimensionar o quadro aumenta. O processo ocupou quase os dois núcleos durante a inferência (\cpuYoloFpSeis\% de CPU no YOLOv8n FP32 a 640) e usou de \rssYoloIntTres{} a \rssYoloFpSeis{}\,MB de memória residente no YOLOv8n, descontados os quadros pré-carregados. Esses resultados sustentam H1 no conjunto medido.

% Gerado por src/analysis/artigo_tex.py. Não edite à mão.
\begin{table}[htbp]
\centering
\footnotesize
\caption{Latência fim a fim por quadro na borda (s; média e IC 95\%), ambiente emulado (QEMU, perfil rpi4b).}
\label{tab:latencia_borda}
\begin{tabular}{llccr}
\hline
Modelo & Resolução & FP32 & INT8 & FP32/INT8 \\
\hline
YOLOv8n & 320 & -- & 4,00 [3,95; 4,05] & -- \\
YOLOv8n & 480 & -- & 8,14 [8,04; 8,23] & -- \\
\hline
\end{tabular}
\par\smallskip\parbox{0.95\linewidth}{\scriptsize Tempos medidos em emulação: servem para comparação relativa, não representam um Raspberry Pi. 300 quadros válidos por combinação.}
\end{table}

% Gerado por src/analysis/artigo_tex.py a partir de results/. Não edite à mão.
\begin{table*}[t]
\centering
\footnotesize
\caption{Baseline de nuvem: latência média por quadro, tamanho do quadro JPEG e banda de subida demandada por taxa da câmera.}
\label{tab:nuvem}
\setlength{\tabcolsep}{4pt}
\begin{tabular}{@{}llrrrr*{3}{r}@{}}
\toprule
& & \multicolumn{3}{c}{Latência média (ms)} & & \multicolumn{3}{c}{Banda de subida (Mbit/s)} \\ \cmidrule(lr){3-5} \cmidrule(l){7-9}
Modelo & {Res. (px)} & {fim a fim} & {transm.} & {infer.} & {kB/quadro} & {5 FPS} & {15 FPS} & {30 FPS} \\
\midrule
YOLOv8n & 320 & 77 & 38 & 16 & 20,7 & 0,8 & 2,5 & 5,0 \\
 & 480 & 111 & 57 & 27 & 43,6 & 1,7 & 5,2 & 10,5 \\
 & 640 & 165 & 82 & 52 & 73,9 & 3,0 & 8,9 & 17,7 \\
SSD-MobileNet V1 & 320 & 87 & 38 & 19 & 20,7 & 0,8 & 2,5 & 5,0 \\
 & 480 & 102 & 57 & 13 & 43,6 & 1,7 & 5,2 & 10,5 \\
 & 640 & 136 & 83 & 17 & 73,9 & 3,0 & 8,9 & 17,7 \\
EfficientDet-Lite0 & 320 & 91 & 38 & 25 & 20,7 & 0,8 & 2,5 & 5,0 \\
 & 480 & 106 & 57 & 17 & 43,6 & 1,7 & 5,2 & 10,5 \\
 & 640 & 137 & 82 & 19 & 73,9 & 3,0 & 8,9 & 17,7 \\
\bottomrule
\end{tabular}
\par\smallskip
\parbox{0.96\linewidth}{\scriptsize Servidor x86-64 nativo, modelos FP32; enlace emulado com tc netem (RTT 40 ms, subida 10 Mbit/s). A banda demandada independe do enlace. Na borda nenhum quadro de vídeo é enviado.}
\end{table*}


Nos cenários derivados (Figura~\ref{fig:descarte}), o YOLOv8n FP32 a 640 pixels sustenta \fpsYoloFpSeis{} quadros por segundo: a 15 e a 30 FPS a fila cresce continuamente e uma política de descarte perderia \descYoloFpSeisQuinze\% e \descYoloFpSeisTrinta\% dos quadros. Em INT8, todas as combinações sustentam 30 FPS, exceto o YOLOv8n a 640 pixels (\fpsYoloIntSeis{} FPS). Em FP32, SSD-MobileNet~V1 e EfficientDet-Lite0 sustentam 15 FPS, mas não 30.

\begin{figure}[!ht]
\centering
\includegraphics[width=\columnwidth]{e6_borda_descarte.pdf}
\caption{Quadros que seriam descartados na borda com a política de processar o quadro mais recente (cenários derivados, ARM64 nativo com 2 núcleos).}
\label{fig:descarte}
\end{figure}

\subsection{Acelerador Edge TPU}

Sem o acelerador físico, avaliou-se a compilação (Tabela~\ref{tab:edgetpu}). Só o SSD-MobileNet~V1 INT8 compilou, com todas as operações mapeadas para a TPU e o modelo inteiro cabendo na memória do chip. O YOLOv8n INT8 foi recusado por não ser inteiramente quantizado e o EfficientDet-Lite0 INT8 falhou em um operador de soma. Sem mudanças nos modelos, só o SSD-MobileNet~V1 aproveitaria o acelerador.

% Gerado por src/analysis/artigo_tex.py. Não edite à mão.
\begin{table}[htbp]
\centering
\footnotesize
\caption{Compilação dos modelos INT8 para o Edge TPU (edgetpu\_compiler 14.1).}
\label{tab:edgetpu}
\begin{tabular}{llcrrl}
\hline
Modelo & Res. & Compilou & Op. na TPU & Op. total & Memória no chip \\
\hline
EfficientDet-Lite0 & 320 & não & -- & -- & -- \\
SSD-MobileNet V1 & 300 & sim & 64 & 64 & 6.99MiB \\
YOLOv8n & 320 & não & -- & -- & -- \\
YOLOv8n & 480 & não & -- & -- & -- \\
YOLOv8n & 640 & não & -- & -- & -- \\
\hline
\end{tabular}
\end{table}


\subsection{Borda $\times$ nuvem (QP2)}

No baseline de nuvem (Tabela~\ref{tab:nuvem}), cada quadro JPEG exige banda de subida que cresce com a resolução e com a taxa da câmera: a 15 FPS são \bandaTresQuinze{}, \bandaQuatroQuinze{} e \bandaSeisQuinze{}\,Mbit/s a 320, 480 e 640 pixels, e a 30 FPS o quadro de 640 pixels exigiria \bandaSeisTrinta{}\,Mbit/s, acima do enlace de subida emulado de 10\,Mbit/s (Figura~\ref{fig:banda}). Na borda nenhum quadro de vídeo sai do dispositivo, apenas o resultado da detecção, o que sustenta H2. A latência fim a fim da nuvem (de \etoeYoloFpTresNuvem{} a \etoeYoloFpSeisNuvem{}\,ms no YOLOv8n) inclui a transmissão do quadro, que cresce com a resolução. Ela não é comparada em valor absoluto com a da borda, porque os dois lados rodaram em servidores diferentes e nenhum deles é um Raspberry Pi.

\begin{figure}[!ht]
\centering
\includegraphics[width=\columnwidth]{e6_banda_subida.pdf}
\caption{Banda de subida necessária para enviar o vídeo à nuvem (JPEG qualidade 80), por taxa e resolução da câmera, comparada à borda, que não envia vídeo.}
\label{fig:banda}
\end{figure}

\subsection{Discussão}

Os resultados indicam três pontos. Primeiro, a quantização INT8 foi o fator que mais reduziu a latência medida, com custo pequeno de qualidade, o que está de acordo com o que a literatura relata \pendente{citar T2}. O ganho variou entre os modelos: foi maior no YOLOv8n, mesmo com o INT8 híbrido, e menor no EfficientDet-Lite0. Segundo, há uma troca entre qualidade e custo que depende da resolução apenas no YOLOv8n: subir de 320 para 640 pixels aumentou o mAP@50 e quase quadruplicou a latência, enquanto nos modelos de entrada fixa a resolução da câmera pouco mudou qualquer das duas medidas. Terceiro, a banda de subida da nuvem cresce linearmente com a taxa da câmera e, nas resoluções e taxas mais altas, supera o enlace residencial assumido; esse custo não existe na borda. Quanto ao Edge TPU, só o SSD-MobileNet~V1 compilou sem alterações, o que vai na mesma direção do que a literatura relata para o YOLOv8n no acelerador \pendente{citar T1}.

\subsection{Limitações}

Os tempos de borda vêm de um núcleo ARM de servidor (Neoverse-N1), mais rápido que o Cortex-A72 do Raspberry Pi~4, e de 2 núcleos em vez de 4; servem para comparar modelos, precisões e resoluções em ARM64 real, não como tempo absoluto de uma placa. Pelo mesmo motivo, a latência da borda não foi comparada em valor absoluto com a da nuvem. A rede da nuvem usa valores assumidos de atraso e banda. A amostragem foi reduzida a uma medição de 300 quadros válidos por combinação, e os cenários de taxa da câmera foram derivados dela, não medidos em execução própria. As imagens são do COCO, sem gravações residenciais. A energia não foi calculada e o Edge TPU foi avaliado apenas pelo relatório do compilador.
